\section{Deployment Scenarios and Validation Plan}
\label{sec_case}
Camera and inference pipelines are motivating scenarios rather than case studies.
A camera study must identify resolution, frame rate, HDR burst length, segmentation
model, tensor dimensions, correctness criterion, and background state.  An
inference or AR study must identify model, scene, tracking-quality metric, reload
behavior, and task-completion criterion.  Each must publish repetitions and
variance rather than phrases such as ``no noticeable delay.''

An Android prototype must disclose the AOSP tag, ART and kernel commits, GKI/vendor
changes, collector, heap limit, zRAM, LMKD properties, MGLRU state, root/bootloader
requirements, changed files and lines of code, service permissions, communication
path, and SELinux policy.  If a custom runtime hook is introduced, its public
contract and synchronization cost must be included; the name ``GCGuard API'' is
not used because no such implementation exists in this repository.

Compatibility must be stated narrowly as the tested application set, not as
compatibility with all Android software.  Energy and sustainability claims require
CPU, DRAM, device-power, wakeup, compression-time, thermal, and battery
measurements.  Until these experiments and artifacts exist, HMS remains a
simulation-backed design proposal.
